\documentclass[xcolor={usenames,x11names},compress]{beamer}
% Otherwise use xcolor=x11names

%% General document %%%%%%%%%%%%%%%%%%%%%%%%%%%%%%%%%%
\usepackage{graphicx}
\graphicspath{{graphics/}}

\usepackage{tikz}
% \usetikzlibrary{decorations.pathreplacing}
% \tikzstyle{every picture}+=[remember picture]

\usepackage{array}
\newcolumntype{P}[1]{>{\raggedright\arraybackslash}p{#1}}

\usepackage{color,colortbl}
\input{inputs/rgb}

% \usepackage{animate}

% \usepackage[table]{xcolor}
  \usepackage{multirow}

\definecolor{CMEEBlue}      {RGB}{99,0, 230}
\definecolor{CMEELightBlue} {RGB}{30,144,255}
\definecolor{CMEEVLightBlue}{RGB}{135,206,250}

%%%%%%%%%%%%%%%%%%%%%%%%%%%%%%%%%%%%%%%%%%%%%%%%%%%%%%

%% Beamer Layout %%%%%%%%%%%%%%%%%%%%%%%%%%%%%%%%%%

\usetheme{Madrid}
\setbeamertemplate{headline}{}
\title{QMEE/CMEE Intro}
\author{Samraat Pawar \& James Rosindell}

\setbeamertemplate{footline}
{
  \leavevmode%
  \hbox{%
%   \begin{beamercolorbox}[wd=.4\paperwidth,ht=2.25ex,dp=1ex,center]{author in 
% head/foot}%
%     \usebeamerfont{author in head/foot}\insertshortauthor
%   \end{beamercolorbox}%
  \begin{beamercolorbox}[wd=.93\paperwidth,ht=2.25ex,dp=1ex,left]{title in 
head/foot}%
    \usebeamerfont{title in head/foot}\insertshorttitle\hspace*{3em}
  \end{beamercolorbox}}%
  \begin{beamercolorbox}[wd=.07\paperwidth,ht=2.25ex,dp=1ex,center]{}
     \insertframenumber{} / \inserttotalframenumber\hspace*{1ex}
%       \insertframenumber{}
  \end{beamercolorbox}

%   \vskip0pt%
}

% \useoutertheme[subsection=false,shadow]{miniframes}
\useinnertheme{default}
\usefonttheme{serif}
\usepackage{palatino}

\setbeamerfont{title like}{shape=\scshape}
\setbeamerfont{frametitle}{shape=\scshape, series = \bfseries}
\setbeamertemplate{frametitle}[default][center]
\setbeamertemplate{note page}[plain]

\setbeamercolor*{lower separation line head}{bg=DeepSkyBlue4} 
\setbeamercolor*{normal text}{fg=black,bg=white} 
\setbeamercolor*{alerted text}{fg=red} 
\setbeamercolor*{example text}{fg=black} 
\setbeamercolor*{structure}{fg=black}
 
\setbeamercolor*{palette tertiary}{fg=black,bg=black!10} 
\setbeamercolor*{palette quaternary}{fg=black,bg=black!10} 

\renewcommand{\(}{\begin{columns}}
\renewcommand{\)}{\end{columns}}
\newcommand{\<}[1]{\begin{column}{#1}}
\renewcommand{\>}{\end{column}}

\def\signed #1{{\leavevmode\unskip\nobreak\hfil\penalty50\hskip2em
  \hbox{}\nobreak\hfil(#1)%
  \parfillskip=0pt \finalhyphendemerits=0 \endgraf}}

\newsavebox\mybox
\newenvironment{aquote}[1]
  {\savebox\mybox{#1}\begin{quote}}
  {\signed{\usebox\mybox}\end{quote}}
  
%%%%%%%%%%%%%%%%%%%%%%%%%%%%%%%%%%%%%%%%%%%%%%%%%%%%

\begin{document}

%%%%%%%%%%%%%%%%%%%%%%%%%%%%%%%%%%%%%%%%%%%%%%%%%%%%%%
\begin{frame}[plain]

\title{Quantitative and Computational Methods in Ecology \& Evolution:
Introduction}
\vspace{12pt}
% \subtitle{}
\author{
	Samraat Pawar \& James Rosindell\\
	\vspace{10pt}
	{\it Department of Life Sciences\\
	\vspace{10pt}
	 Silwood Park}\\
	\vspace{20pt}
  \centering
  \includegraphics[height = .3in]{Imperial_Color1.pdf}
}
 
\titlepage
\end{frame}

%%%%%%%%%%%%%%%%%%%%%%%%%%%%%%%%%%%%%%%%%%%%%%%%%%%%%%
% \section{\scshape Introduction}
% \subsection{}
%%%%%%%%%%%%%%%%%%%%%%%%%%%%%%%%%%%%%%%%%%%%%%%%%%%%%%

\begin{frame}{Why Ecology and Evolution?}

\centering

\includegraphics[width=.7\textwidth]{Pieter_Fish.jpg}

\vspace{6pt}

{\small {\bf \it Big Fish Eat Little Fish}, 1557, Pieter van der Heyden}

\end{frame}

%%%%%%%%%%%%%%%%%%%%%%%%%%%%%%%%%%%%%%%%%%%%%%%%%%%%%%
\begin{frame}{Why Ecology and Evolution?}

There is a long list of real world applications that need Quantitative innovation and Computational methods in Ecology and Evolution ...
 
  \begin{itemize}[<+->] \setlength{\itemindent}{0em}\itemsep6pt

		\item Agro-ecology and sustainable agricuture 
		
		\item Economic systems (micro-economics in particular) 

		\item Disease emergence and transmission/outbreak (especially complex transmission systems such as vectro-borne diseases)
		 
	 \item Global climate change, carbon balance, nitrogen cycling
	 
	 \item Ecosystem engineering (think yogurt cheese \& wine!)
		 
	 \item Species invasions and ecosystem collapse

  \end{itemize} 

    % \pause
% \centering
    % {\bf ergo, Computing!}\\
  % \includegraphics[width=.4\textwidth]{IBM_Blue_Gene.jpg}

\end{frame}

%%%%%%%%%%%%%%%%%%%%%%%%%%%%%%%%%%%%%%%%%%%%%%%%%%%%%%

\begin{frame}{We live in ``interesting'' times}

\pause
\begin{columns}[c]
  \column{0.6\linewidth}
		\begin{itemize}[<+->]\setlength{\itemindent}{0em}
			\item Climatic warming and fluctuations
		\end{itemize}
  \column{0.4\linewidth}\centering

		% {\tiny 5-yr mean surface temperature anomaly (NASA 2013)}
			\includegraphics[width=.9\textwidth]<1>{Composite/001.jpg}
			\includegraphics[width=.9\textwidth]<2>{Composite/148.jpg}
			\includegraphics[width=.9\textwidth]<3>{Composite/295.jpg}
			\includegraphics[width=.9\textwidth]<4>{Composite/443.jpg}
			\includegraphics[width=.9\textwidth]<5>{Composite/590.jpg}
			\includegraphics[width=.9\textwidth]<6>{Composite/800.jpg}		

\end{columns}

\pause

\begin{columns}[c]
  \column{0.6\linewidth}
		\begin{itemize}[<+->]\setlength{\itemindent}{0em}
			\item Species range shifts and invasions 
		\end{itemize}
  \column{0.4\linewidth} \centering
		\includegraphics[width=.9\textwidth]{Kudzu.pdf}		
\end{columns}

\pause

\begin{columns}[c]
  \column{0.6\linewidth}
		\begin{itemize}[<+->]\setlength{\itemindent}{0em}
			\item Overexploitation of ecosystems
		\end{itemize}
  \column{0.4\linewidth}\centering
		\includegraphics[width=.9\textwidth]{Pieter_Fish.jpg}		
\end{columns}

\end{frame}

%%%%%%%%%%%%%%%%%%%%%%%%%%%%%%%%%%%%%%%%%%%%%%%%%%%%%%

\begin{frame}{Some worrying effects on (complex) biological systems}

\begin{columns}[c]
  \column{0.55\linewidth}
		\begin{itemize}[<+->]\setlength{\itemindent}{0em}
			\item Loss or collapse of ecosystem function
			
			\item Loss of ecosystem recovery or succession

		\end{itemize}

  \column{0.45\linewidth}
    \begin{figure}
			\includegraphics[width=\textwidth]{Paleo.pdf}		
		\end{figure}
\end{columns}

\pause

\begin{columns}[c]
  \column{0.55\linewidth}
		\begin{itemize}\setlength{\itemindent}{0em}
			\item Disease emergence \& outbreak
		\end{itemize}
  \column{0.45\linewidth}

	% \begin{center}
		\begin{tikzpicture}
		  \node (img1) {\includegraphics[width=\textwidth]{Chagas.pdf}};
		  \pause
		  \node (img2) at (img1) {\includegraphics[width=\textwidth]{Chagas_box.pdf}};
		\end{tikzpicture}
	% \end{center}
	
  % \begin{figure}
		% \includegraphics[width=\textwidth]<1>{Chagas.pdf}		
		% \includegraphics[width=\textwidth]<2>{Chagas_box.pdf}
	% \end{figure}
		
\end{columns}

\end{frame}

%%%%%%%%%%%%%%%%%%%%%%%%%%%%%%%%%%%%%%%%%%%%%%%%%%%%%%

\begin{frame}{Why Ecology and Evolution?}
\centering
	\includegraphics[width=.55\textwidth]{EcoEvoRoadmap.pdf}
\end{frame}

%%%%%%%%%%%%%%%%%%%%%%%%%%%%%%%%%%%%%%%%%%%%%%%%%%%%%%
\begin{frame}{Why {\it Computational} Ecology and Evolution?}

  \begin{itemize}[<+->] \setlength{\itemindent}{0em}\itemsep3pt

    \item Biological systems are complex (e.g., nonlinear dynamics are rife)

    \item Maths is far more powerful with numerical analyses

		\item Biological data are BIG and MESSY --- need to be organized 
		and managed 

		\item Biological data need MINING and often, complex STATISTICAL ANALYSES

		\item Biology needs REPRODUCIBLE analyses and work-flows 
		that link theory and data \pause {\it -- employers want 
		this too!}

  \end{itemize} 

\centering
    % {\bf ergo, Computing!}\\
  \includegraphics[width=.4\textwidth]{IBM_Blue_Gene.jpg}

\end{frame}

%%%%%%%%%%%%%%%%%%%%%%%%%%%%%%%%%%%%%%%%%%%%%%%%%%%%%%
\begin{frame}[plain]{}

\begin{aquote}{Donald Knuth}
It is hard for me to say confidently that, after fifty more years of explosive 
growth of computer science, there will still be a lot of fascinating unsolved 
problems at peoples' fingertips, that it won't be pretty much working on 
refinements of well-explored things. Maybe all of the simple stuff and the 
really great stuff has been discovered. It may not be true, but I can't predict 
an unending growth. I can't be as confident about computer science as I can 
about biology. Biology easily has 500 years of exciting problems to work on, 
it's at that level.
\end{aquote}
\end{frame}


%%%%%%%%%%%%%%%%%%%%%%%%%%%%%%%%%%%%%%%%%%%%%%%%%%%%%%
\begin{frame}{Why {\it Computational} Ecology and Evolution?}

\centering

\includegraphics[width=.9\textwidth]{Joel_Microscope.png}

  \vspace{20pt}

  \pause 

  {\bf Read it}  --- its on the course's Git repository \\\url{https://github.com/mhasoba/TheMulQuaBio}

\end{frame}

%%%%%%%%%%%%%%%%%%%%%%%%%%%%%%%%%%%%%%%%%%%%%%%%%%%%%%
\begin{frame}{Computational Ecology and Evolution at {\it Silwood}}

{\bf A wide range of theoretical and empirical research (big data!)}:
\scriptsize
  \begin{itemize}\setlength{\itemindent}{0em}\itemsep2pt
		\item Evolution and Developmental Genetics (Abzhanov)
    \item Genetics and behavior (Schroeder)
    \item Tropical biology (Ewers, Banks-Leite)
    \item Vector borne diseases (Cator, Burt)
    \item Phylogenetics, genomics (Savolainen, Fumagalli)
    \item Human genetics (Hodgson)
    \item Behavior (Cator, Gill)
    \item Paleontology (Brazeau)
    \item Food webs, networks (Woodward, O'Gorman, Pawar)
    \item Metabolic theory, population biology (Pawar) 
    \item Pollinator behavior and ecology (Gill)
    \item Neutral theory, scientific visualization (Rosindell)
    \item Population genetics (Burt, Schroeder, Hodgson, Barraclough, Fumagalli)
    \item Conservation biology (Knight, Mills, Banks-Leite, Ewers)
    \item Microbial systems (Bell, Barraclough)
    \pause
    \item {\bf Many more PIs across are multiple campuses!}
  \end{itemize}  

\end{frame}

%%%%%%%%%%%%%%%%%%%%%%%%%%%%%%%%%%%%%%%%%%%%%%%%%%%%%%
\begin{frame}{What you will learn}
  
  \begin{itemize}\setlength{\itemindent}{0em}\itemsep6pt
\item Competence in stat-of-art quantitative methods for addressing 
	modern biological problems
	
	\item How to select the correct quantitative tool to address a 
	specific biological problem

	\item An ability to develop, analyse, numerically simulate, fit 
	models to data and interpret

	\item Quantitative models of biological systems, including 
	statistical and mathematical models

 \item Techniques in Population biology, Population genetics, Genomics 
 and Geographical Information Systems

 \item How to design and conduct research, with the necessary 
 computational workflows -- {\it please have a look at past projects!}
 
 \end{itemize}  

  % {\it Note that MSc's have about double the coursework in these areas}
  
\end{frame}


%%%%%%%%%%%%%%%%%%%%%%%%%%%%%%%%%%%%%%%%%%%%%%%%%%%%%%
\begin{frame}{Course Administration}

\small
\begin{center}
  \begin{tabular}{>{\it}lP{0.52\textwidth} lP{0.6\textwidth}}
\rowcolor{CMEEVLightBlue}    
    Course Director & Dr. Samraat Pawar (ext. 42213, 
  s.pawar@imperial.ac.uk)\\
    Course Co-Director  & Dr. James Rosindell (ext. 42242, j.rosindell@imperial.ac.uk)\\
\rowcolor{CMEEVLightBlue}    
    Postgraduate Administrator	  & Mrs. Amanda Ellis (ext. 42251,  amanda.ellis@imperial.ac.uk)\\
    Postgraduate Tutor	& Dr. Julia Schroeder (julia.schroeder@imperial.ac.uk)\\
\rowcolor{CMEEVLightBlue}    
		Director of Postgraduate Studies &   Dr. Niki Gounaris (ext. 4 5209, k.gounaris@imperial.ac.uk)\\
		Course Tutors		& Katie Willis (katie.willis16@imperial.ac.uk) \& Tom Clegg (t.clegg17@imperial.ac.uk) \\
\rowcolor{CMEEVLightBlue}    
    Course Representative	& Up to you (see Silwood Masters Guidebook)\\
  \end{tabular}

\end{center}
{\it Add 020 759 to extension numbers to call from external phones}

\end{frame}

%%%%%%%%%%%%%%%%%%%%%%%%%%%%%%%%%%%%%%%%%%%%%%%%%%%%%%
\begin{frame}{Getting started}
  
  \begin{itemize}\setlength{\itemindent}{0em}\itemsep6pt
    \item You each will receive a laptop computer and peripherals -- {\it If you want it}  
    \item You will be provided a USB stick with Ubuntu 16.04 64 bit -- {\it If you need it (ubuntu 18.04 will also work)}
    \item If you are using your own laptop -- use Ubuntu/Lubuntu/Xubuntu 14.04 64 bit or higher, or Mac OS
    \begin{itemize}
	    \item If using Mac OS, remember that not all commands behave the same between Mac OS and Linux 
		\end{itemize}
    
		\item Make sure you can access secured Imperial College wireless (use college name and password) --- external Bootcampers, use Eduroam or a guest account (ask me later) 
    \item Get an bitbucket or github account (read their Git tutorials (quite intuitive)) --- free private repos using your college email account
     \item You will be assisted by very capable demonstrators -- also, learn collaboratively with your your classmates!
  
\end{itemize}  
 
  \end{frame}

  %%%%%%%%%%%%%%%%%%%%%%%%%%%%%%%%%%%%%%%%%%%%%%%%%%%%%%
\begin{frame}{Computer guidelines and rules}
  
  \begin{itemize}\setlength{\itemindent}{0em}\itemsep6pt
    \item You are responsible for your computer hardware and software (we have your money!)
    \item You should be able to install all necessary (open source) software
    \item Your computer is your main tool, you are expected to achieve a high
degree of mastery of it!
    \item Return it to Martin Selby at end of Course (CMEE Masters) or Bootcamp/Winter term (QMEE CDT)
    \item We expect it to be undamaged
  \item Please do not leave your computer in any room (other than your
residence!) overnight, starting now
 \end{itemize}  
  \end{frame}

    %%%%%%%%%%%%%%%%%%%%%%%%%%%%%%%%%%%%%%%%%%%%%%%%%%%%%%
\begin{frame}{Handbook and Lectures}
  
  \begin{itemize}\setlength{\itemindent}{0em}\itemsep6pt
%     \item Printed guidebooks may become outdated, download updated ones from
% the git repo
    \item Please check key dates for coursework and reports in your gudebook
    \item Lectures:
    \begin{itemize}\setlength{\itemindent}{-1em}
      \item 2 1-hr lectures in the morning (1000 -- 1230, except in come cases)
      \item 3-hour practical in afternoon, except on Wednesdays
      \item Lecturers will stay for at least 1 hour of practical session
      \item There may be deviations from this -- check updated timetables (use ical)!
      \item Demonstrators will be available during practicals in botcamp and most other weeks
      \item Most lectures in CPB or Hamilton Comp room in first 9 weeks, then Wallace Room after MRes cohort finishes
    \end{itemize}
 \end{itemize}  
  \end{frame}

%%%%%%%%%%%%%%%%%%%%%%%%%%%%%%%%%%%%%%%%%%%%%%%%%%%%%%
\begin{frame}{Seminars}

  \begin{itemize}\setlength{\itemindent}{0em}\itemsep12pt

    \item Thursday seminars at 1300 hrs in this building (web link in guidebook)
    \item Students must attend Thursday seminars -- 1/2 page summary each 
of min 12 seminars (seminar diary) due after seminar series ends in summer
	 \item You are encouraged to give talks for feedback / discussion at 
	 other times of the week, especially in the Spring --- great for 
	 running ideas past peers (you can give multiple short ones)

	 \item There bi-weekly Social Seminar Series (Beer, Pizza, Talks!) is also a great platform to give a talk.   
	  
	  \end{itemize}

  \end{frame}

  %%%%%%%%%%%%%%%%%%%%%%%%%%%%%%%%%%%%%%%%%%%%%%%%%%%%%%
\begin{frame}{Workshops and Symposia}

  \begin{itemize}\setlength{\itemindent}{0em}\itemsep12pt

    \item Workshops organized by us are all optional, but strongly recommended -- check guidebook
    \item Some important ones this week

    \item Summer graduate symposium on Frontiers in Ecology and Evolution (FrEEC Symposium -- \url{freesymp.org}):
    \begin{itemize}\setlength{\itemindent}{-1em}
      \item Week 2 of September, organized by Masters + PhD students + Postdocs
      \item More details in Silwood Masters Student Guidebook
    \end{itemize}
 \end{itemize}  
  \end{frame}

%%%%%%%%%%%%%%%%%%%%%%%%%%%%%%%%%%%%%%%%%%%%%%%%%%%%%%
\begin{frame}{Assessment and Marking}

\small

\begin{center}
  \begin{tabular}{>{\it}P{3cm}|P{3cm}|P{4.9cm}}
    \hline
    \rowcolor{CMEEVLightBlue}{\bf Activity} &  {\bf MSc CMEE} & {\bf MRes 
CMEE}\\ \hline
      Lectures + practicals, with assessment & Required for 19 weeks &
Required for first 9 weeks, optional attendance in MSc modules within reason
thereafter \\ \hline
    Exams & Required & Not required \\
\hline
  Project report (Dissertation) & Required & Required \\ \hline
   Seminars & Required, seminar diary to be submitted based on
the attendance of a minimum of 12 seminars & Required, seminar diary to be submitted based on
the attendance of a minimum of 12 seminars\\ \hline
  Workshops & All optional & All optional \\
\hline

\end{tabular}
  
\end{center}
  \end{frame}

%%%%%%%%%%%%%%%%%%%%%%%%%%%%%%%%%%%%%%%%%%%%%%%%%%%%%%
% \begin{frame}{Assessment and Marking}

% \scriptsize

% \begin{center}

    % \begin{tabular}{>{\it}m{3.5cm}|p{1.5cm}|p{1.5cm}|p{1.5cm}|p{1.5cm}}
    % \hline 
		% \rowcolor{CMEEVLightBlue} &  \multicolumn{2}{c|}{\bf MSc CMEE} & \multicolumn{2}{c}{\bf MRes CMEE}\\ \hline
    
    % \rowcolor{CMEEVLightBlue} {\bf Component} &  {\% of Course} & {\% of Component} & {\% of Course} & {\% of Component}\\ \hline

    % \multicolumn{5}{l}{\bf Coursework} \\ \hline 
    % Computing  &  {13.75} & {55} & {13.75} & {55}\\ \hline
    % CMEE Mini-project  &  {6} & {24} & {6} & {24}\\ \hline
    % HPC Long Practical  &  {5} & {20} & {5} & {20} \\ \hline
    % Seminar Diary  &  {0.25} & {1} & {0.25} & {1}\\ \hline
    % {\bf Coursework Total} &  {25} & {100} & {25} & {100}\\ \hline
    % \multicolumn{5}{l}{} \\[-10pt]\hline 

    % \multicolumn{5}{l}{\bf Exams} \\ \hline 
    % Exam 1  &  {10} & {40} & {--} & {--}\\ \hline
    % Exam 2  &  {15} & {60} & {--} & {--}\\ \hline
		% {\bf Exam Total} &  {25} & {100} & {--} & {--}\\ \hline    
    % \multicolumn{5}{l}{} \\[-10pt]\hline 
    
    % \multicolumn{5}{l}{\bf Project} \\ \hline 
    % Final Report + Presentation  &  {35} & {70} & {52.5} & {70}\\ \hline
    % Viva  &  {12.5} & {25} & {18.75} & {25}\\ \hline
    % Supervisor mark  &  {2.5} & {5} & {3.75} & {5}\\ \hline
		% {\bf Project Total} &  {50} & {100} & {75} & {100}\\ \hline    

    % \hline
  % \end{tabular}
% \end{center}

  % \end{frame}

%%%%%%%%%%%%%%%%%%%%%%%%%%%%%%%%%%%%%%%%%%%%%%%%%%%%%%
% \begin{frame}{Important dates}

% \begin{center}
	    %% \rowcolors{1}{CMEEVLightBlue}{}
  % \begin{tabular}{p{3cm}|p{7cm}}
    % \hline
      % \rowcolor{CMEEVLightBlue}{\bf Date} & {\bf Activity/Item due} \\ 
      % \noalign{\smallskip}
	% 9 Dec, 5PM 	& MRes: Project proposal\\ 
	% \rowcolor{CMEEVLightBlue}
	% 15 Dec, 5PM 	& MSc, MRes: HPC Long Practical\\ 
	% 17 Feb, 5PM 	& MSc, MRes: CMEE Miniproject\\ 
	% \rowcolor{CMEEVLightBlue}
	% 7 April, 5PM 	& MSc, Project proposal\\ 
	% 25 Aug, 5pm 	& MSc, MRes: Seminar Diary\\ 
   % \hline 
  % \end{tabular}\\
% \end{center}

% \pause {\centering Further important dates, including thesis submission are same across 
% all Silwood Masters courses -- \bf please refer to the Silwood Student 
% Guidebook}

% \end{frame}

%%%%%%%%%%%%%%%%%%%%%%%%%%%%%%%%%%%%%%%%%%%%%%%%%%%%%%
\begin{frame}{Warm-up for the Bootcamp and rest of the course}

\begin{itemize}\itemsep8pt

\item Get Linux on your computer (if applicable)! (google - ``Install Ubuntu'')
\item Lots of UNIX tutorials out there. Try 
\url{https://software-carpentry.org/lessons/} (Chapter ``shell''). 
\item Excellent book on Git: \url{http://git-scm.com/book}, also,
\url{https://www.atlassian.com/git/}
\item See \url{http://www.andy-roberts.net/writing/latex/benefits}
\item Also, Word vs. \LaTeX: \url{http://openwetware.org/wiki/Word_vs._LaTeX}

\end{itemize}

\centering
{\it More extensive list in Guidebook and Course notes (are you getting my emails?)}

\end{frame}

%%%%%%%%%%%%%%%%%%%%%%%%%%%%%%%%%%%%%%%%%%%%%%%%%%%%%%
\begin{frame}{Rest of the Course/Bootcamp}

\begin{itemize}\itemsep12pt

\item A number of Readings/Resources on the git repo 

\item Read, Read, Read about concepts and models in Ecology \& Evolution!

\item Happy QMEE (Boot)campers --- if you plant to attend any of the weeks beyond the Bootcamp, please let Joel Santana-Mondo know!
		\begin{itemize}
    \item You may skip parts of the Bootcamp that are too basic for you 
    \item You need to return your Laptop (if applicable) latest by end of Winter term 
		\end{itemize}
		
		
\end{itemize}

\end{frame}
%%%%%%%%%%%%%%%%%%%%%%%%%%%%%%%%%%%%%%%%%%%%%%%%%%%%%%
\begin{frame}{Questions?}

\centering
  \includegraphics[width=\textwidth]{CMEE_Fungus.jpg}\\
    {\it \small (soon to be) Famous CMEE Fungus}
  
\end{frame}



\end{document}
